\section{Kimi K2: Open Agentic Intelligence}

Las dos secciones anteriores trataron el marco general de los Agent; esta sección aborda el primer informe técnico, y su pregunta central es cómo un modelo abierto puede entrenarse de forma sistemática para adquirir capacidades Agentic y acercarse a la frontera de los modelos cerrados. Kimi K2 es un modelo MoE del orden de 1T de parámetros, con unos 32B de parámetros activos. El informe propone el optimizador MuonClip, que usa QK-Clip para resolver la inestabilidad del entrenamiento con Muon sin perder la token efficiency; K2 se preentrenó con 15.5T tokens, y el posentrenamiento incluye agentic data synthesis a gran escala y joint RL.

\begin{figure}[H]
\centering
\includegraphics[width=0.96\textwidth]{kimi-k2-agentic-pipeline.png}
\caption{Kimi K2 Agentic Pipeline: MuonClip, 15.5T tokens, agentic data synthesis y joint RL. Diagrama conceptual de elaboración propia, basado en el informe técnico de Kimi K2 y en la conversación de 00:30:57--00:43:50.}
\end{figure}

\begin{knowledgebox}{Lectura de la figura: lo importante de Kimi K2 no es un solo benchmark}
K2 reúne la base MoE, un optimizador estable, datos a gran escala, la síntesis de trayectorias de herramientas y RL. Su capacidad agentic proviene de un pipeline de entrenamiento completo, no solo de un procesamiento posterior.
\end{knowledgebox}

\subsection{MuonClip y QK-Clip}

Para entender K2 no basta con ver los agentic benchmark; también hay que ver cómo logra un entrenamiento estable. Esta subsección explica MuonClip y QK-Clip. Muon tiene la ventaja de la token efficiency, pero el entrenamiento a gran escala puede volverse inestable; QK-Clip controla el riesgo de divergencia del entrenamiento al acotar los attention logits. El informe subraya que MuonClip permitió que K2 se mantuviera estable durante el entrenamiento con 15.5T tokens.

\begin{knowledgebox}{Términos clave: componentes del entrenamiento de Kimi K2}
\begin{tabularx}{\textwidth}{p{0.22\textwidth}p{0.34\textwidth}X}
\toprule
Componente & Función & Riesgo/importancia \\
\midrule
MoE & Aumenta el total de parámetros, pero controla el cómputo activo & El enrutamiento, la carga de los expertos y la estabilidad del entrenamiento son complejos. \\
MuonClip & Mejora la token efficiency y estabiliza el entrenamiento & Debe manejar la explosión de los attention logits. \\
QK-Clip & Acota los logits que producen las proyecciones de query/key & Es un mecanismo de estabilización de intervención mínima. \\
Agentic data synthesis & Genera trayectorias de uso de herramientas y datos de tareas & El filtrado de calidad y el diseño de la recompensa son decisivos. \\
Joint RL & Mejora las capacidades en entornos reales o sintéticos & Los entornos y la reward tienen un costo alto. \\
\bottomrule
\end{tabularx}
\end{knowledgebox}

\subsection{Resumen de la sección}

El valor didáctico de Kimi K2 es que combina el modelo abierto, MoE, la estabilidad del optimizador, las trayectorias sintéticas de herramientas y RL en un agentic training pipeline.
